\documentclass[11pt]{article}
\usepackage{fontspec}
\setmainfont{Times New Roman}
\setsansfont{Arial}
\setmonofont{Consolas}
\usepackage{amsmath,amssymb,mathtools,array}
\usepackage{geometry}
\usepackage{microtype}
\geometry{a4paper,margin=2.35cm}
\setlength{\parindent}{1.25em}
\setlength{\parskip}{0pt}
\makeatletter
\renewcommand\footnotesize{\@setfontsize\footnotesize{7.8}{8.8}\selectfont}
\makeatother
\emergencystretch=35em
\allowdisplaybreaks
\sloppy
\hbadness=10000
\vbadness=10000
\hfuzz=2pt
\vfuzz=10pt
\raggedbottom
\providecommand{\fraka}{\mathfrak a}
\providecommand{\frakb}{\mathfrak b}
\providecommand{\frakp}{\mathfrak p}
\providecommand{\frakq}{\mathfrak q}
\providecommand{\frakR}{\mathfrak R}
\providecommand{\frakS}{\mathfrak S}
\providecommand{\Norm}{N}

\begin{document}
\setcounter{footnote}{0}

% Unit: Paper25 v001. Direct Russian translation from the German final-audited source slice.
\section*{25. Теория элиминации и теория идеалов}

\begin{center}
J. Ber. d. DMV 33 (1924), с. 116--120
\end{center}

В дальнейшем я хотела бы показать, как теорию элиминации можно построить вполне параллельно теории нулей для полиномов от одной переменной; это построение основано на соединении теории элиминации, данной Hentzelt, с методами теории полей, как их развил Steinitz, и теории идеалов, как я сама ее разработала.\footnote{Подробное изложение и библиографические указания: Math. Ann. 90 и работа, которая вскоре появится.}

Сначала я набросаю постановку вопроса в случае полиномов от одной переменной. В качестве области коэффициентов раз и навсегда положим фиксированное поле $P$, к которому относятся понятия простой функции и т. д. При этом $P$ можно считать абстрактно определенным в смысле Steinitz; но, разумеется, в специальном случае оно может совпадать с полем рациональных чисел или с полем всех комплексных чисел --- согласно предположению, ранее общепринятому в теории элиминации. Теория нулей для полиномов от одной переменной теперь распадается на четыре шага:

\noindent 1. \emph{Теорема о разложении}, которая позволяет свести вопрос к вопросу о простых функциях --- полиномах, неприводимых в $P$. Ибо из
\[
 a(x)=p_1(x)^{\varrho_1}\cdots p_r(x)^{\varrho_r}=q_1(x)\cdots q_r(x)
\]
--- где $p(x)$ обозначают различные простые функции, а $q(x)$, следовательно, "наибольшие первичные функции", --- следует, что $a(x)$ обращается в нуль во всех и только в нулях простых функций $p(x)$.

\noindent 2. \emph{Построение поля нулей для простой функции $p(x)$.} Такое поле нулей $\frakR=P(\alpha)$ --- где $\alpha$ обозначает нуль $p(x)$ --- по Steinitz всегда можно построить как поле расширения поля $P$, изоморфное полю классов вычетов $(\frakR)$ по $p(x)$. Обратно, всякое поле нулей, лежащее в поле расширения поля $P$, изоморфно полю классов вычетов $(\frakR)$; поэтому поле нулей и поле классов вычетов абстрактно тождественны.

\noindent 3. \emph{Разложение $p(x)$ на линейные множители в поле Galois, определенном $p(x)$.} Повторяя построение, данное в пункте 2, приходят к существенно однозначно определенному полю Galois $P(\alpha_1,\ldots,\alpha_m)$, которого как раз достаточно, чтобы $p(x)$ разложилось на линейные множители.

\noindent 4. \emph{Значение кратности для исходного полинома $a(x)$.} По пунктам 1, 2, 3 для произвольного полинома $a(x)$ вопрос о нулях и соответствующее разложение на линейные множители решены. Кратностью при этом можно считать степень отдельной первичной функции $q(x)$ --- то есть число линейно независимых относительно $P$ классов вычетов по $q(x)$; ибо в случае алгебраически замкнутой области коэффициентов $P$, где $p(x)$ становятся линейными, это число как раз и дает кратность нулей.

Теперь я перехожу к общей теории элиминации, то есть беру за основу область $\frakS$ всех полиномов от $x_1,\ldots,x_n$ с коэффициентами из $P$; тогда теорию элиминации можно определить как \emph{теорию нулей идеала} из этой полиномиальной области. Под идеалом $\fraka$ я понимаю, как обычно, такую систему полиномов, что вместе с $f$ и $g$ к $\fraka$ принадлежит также $f-g$, а вместе с $f$ также $Af$, где $A$ понимается как произвольный полином из $\frakS$. Под нулем $\fraka$ я понимаю всякую систему значений $\alpha_1,\ldots,\alpha_n$, принадлежащую полю расширения поля $P$, такую, что $f(\alpha)$ обращается в нуль для каждого полинома $f$, принадлежащего идеалу. Если $f_1,\ldots,f_t$ обозначает всегда существующий базис идеала $\fraka$ --- то есть $f=A_1f_1+\cdots+A_tf_t$ для каждого $f$ из $\fraka$ ---, то нули $\fraka$, очевидно, тождественны нулям $f_1,\ldots,f_t$; а так как, обратно, всякую конечную систему полиномов можно рассматривать как базис порожденного ею идеала, то данное выше определение нулей для теории элиминации тождественно обычно употребляемому, но избегает трудности, заключенной в произвольном выделении базиса. Для полиномов от одной переменной, где речь идет только о главных идеалах, эта трудность не возникает, поскольку базисный полином определен однозначно с точностью до единицы --- множителя из $P$; однако и там идеальная точка зрения даст более точное понимание.

Теория элиминации теперь, соответственно случаю одной переменной, распадается на четыре шага:

\noindent I. \emph{Теорема о разложении} --- представление идеала через наибольшие первичные компоненты, --- которая позволяет свести вопрос к вопросу о простых идеалах. Идеал называется простым, если он делит произведение лишь тогда, когда делит по крайней мере один из множителей; первичным, если он делит по крайней мере один из множителей или некоторую степень каждого множителя. Каждому первичному идеалу $\frakq$ принадлежит один и только один ассоциированный простой идеал $\frakp$, который является делителем $\frakq$ и некоторая степень которого делится на $\frakq$. Имеет место представление как наименьшее общее кратное:
\[
 \fraka=[\frakq_1,\ldots,\frakq_r];\qquad
 \frakp_1^{\sigma_1}\cdots \frakp_r^{\sigma_r}\equiv 0\pmod{\fraka}.
\]
Здесь $\frakq$ являются наибольшими первичными компонентами; то есть каждое $\frakq$ первично, но наименьшее общее кратное любых двух $\frakq$ не первично. Кроме того, без ограничения общности можно предполагать, что ни одно $\frakq$ не делит наименьшее общее кратное остальных, так что ни одно $\frakq$ нельзя просто опустить; $\frakp$ являются ассоциированными простыми идеалами соответствующих $\frakq$. Итак, не сами $\frakq$, но --- и в этом состоит аналог случая одной переменной --- их ассоциированные простые идеалы $\frakp$ однозначно определены идеалом $\fraka$; и так как каждое $\frakp$ является делителем $\fraka$, а обратно произведение степеней $\frakp$ делится на $\fraka$, то нули идеала совпадают с нулями его ассоциированных простых идеалов.

\noindent II. \emph{Построение поля нулей для простого идеала $\frakp$.} Классы вычетов по простому идеалу по определению образуют кольцо без делителей нуля, содержащее по крайней мере один элемент, отличный от нулевого, как только $\frakp$ отличается от единичного идеала; образованием частных это кольцо, следовательно, можно расширить до поля классов вычетов $(\frakR)$ идеала $\frakp$. Всегда можно построить поле расширения $\frakR$ поля $P$, изоморфное $(\frakR)$, которое становится полем нулей $\frakp$; то есть $\frakR=P(\alpha_1,\ldots,\alpha_n)$, где $\alpha_1,\ldots,\alpha_n$ обозначает нуль $\frakp$. При этом $\frakR$ имеет степень трансцендентности $k$ $(0\leq k<n)$ относительно $P$, следовательно содержит $k$ алгебраически независимых относительно $P$ элементов, тогда как любые $(k+1)$ элемента становятся алгебраически зависимыми относительно $P$. Обратно, всякое поле нулей степени трансцендентности $k$, лежащее в поле расширения поля $P$, изоморфно полю классов вычетов $(\frakR)$; поэтому поле нулей степени трансцендентности $k$ и поле классов вычетов абстрактно тождественны.

\noindent III. \emph{Разложение нормы $\frakp$ на линейные множители в поле Galois, определенном $\frakp$.} Если считать переменные $x_1,\ldots,x_n$ полученными из исходных переменных $y_1,\ldots,y_n$ линейным преобразованием с неопределенными в качестве коэффициентов, то поле нулей $\frakR$ степени трансцендентности $k$ можно рассматривать как конечное алгебраическое поле расширения поля $P(x_{i+1},\ldots,x_n)$ $(k=n-i)$. Поэтому можно перейти к полю Galois $\frakR$ относительно $P(x_{i+1},\ldots,x_n)$, в котором простая функция $P^{(i)}(x_i,x_{i+1},\ldots,x_n)$ с нулем $\alpha_i$ разлагается на линейные множители. Эта простая функция $P^{(i)}$, если вернуться к исходным переменным $y$, играет роль фундаментального уравнения; ибо $\alpha_i$ становится равным фундаментальной форме $u_1\beta_1+\cdots+u_n\beta_n$, где $\beta$ обозначает систему нулей в переменных $y$, алгебраически зависящую от параметров $x_{i+1},\ldots,x_n$; следовательно, разложение на линейные множители дает произведение по всем нулям. Но простую функцию $P^{(i)}$ или ее степень можно также рассматривать как норму $N(\frakp)$ идеала $\frakp$, если рассматривать $\frakp$ как модуль из линейных форм в степенных произведениях $x_1,\ldots,x_{i-1}$, взять $P(x_{i+1},\ldots,x_n)$ вместо $P$ в качестве области коэффициентов и определить норму как произведение по всем элементарным делителям --- из которых только конечное число отлично от единицы. При совершенном поле $P$ в качестве области коэффициентов всегда имеем $N(\frakp)=P^{(i)}$; при несовершенном поле нормой может оказаться степень $P^{(i)}$.

\noindent IV. \emph{Норма исходного идеала и значение кратности.} Если соответственно последовательно рассматривать идеал $\fraka$ как модуль из линейных форм в степенных произведениях $x_1,\ldots,x_{i-1}$ $(i=1,2,\ldots,n)$, то можно показать, что при взятии за основу $P(x_{i+1},\ldots,x_n)$ этот модуль каждый раз имеет только конечное число элементарных делителей, отличных от единицы; это условие конечности следует из существования базиса идеала. Произведение по этим каждый раз конечным элементарным делителям дает отдельные нормы; норма $N(\fraka)$ определяется как произведение всех отдельных норм. Наибольшие первичные множители $Q^{(i)}(x_i,x_{i+1},\ldots,x_n)$ нормы $N(\fraka)$ взаимно однозначно соответствуют ассоциированным простым идеалам $\frakp$ идеала $\fraka$ так, что $Q^{(i)}$ становится степенью $P^{(i)}$, где $P^{(i)}$ (или, если нужно, степень $P^{(i)}$) обозначает норму $N(\frakp)$ идеала $\frakp$. Тем самым получается разложение
\[
 N(\fraka)=N(\frakp_1)^{\varrho_1}\cdots N(\frakp_r)^{\varrho_r}=Q_1\cdots Q_r,
\]
--- где, если нужно, $N(\frakp)$ следует заменить наивысшим элементарным делителем $E(\frakp)=P^{(i)}$, --- и тем самым по I, II, III для всех нулей $\fraka$ выполнено разложение нормы на линейные множители. Кратностью нулей, соответствующих отдельному простому идеалу $\frakp$ идеала $\fraka$, можно считать степень $Q^{(i)}$, которую можно истолковать как число классов вычетов, линейно независимых относительно $P(x_{i+1},\ldots,x_n)$. А именно речь идет о классах вычетов некоторых изолированных компонент разложения, однозначно определенных $\frakp$ и простыми идеалами большей размерности: о классах вычетов основного идеала $(i-1)$-й ступени по компоненте основного идеала $i$-й ступени, определенной $\frakp$; последняя компонента имеет $\frakp$ и все простые идеалы большей размерности как ассоциированные простые идеалы, первая же имеет только простые идеалы большей размерности.

Полное \emph{включение случая одной переменной} получается, если и здесь перейти к главному идеалу. Разложение, данное в пункте 1, тогда --- смотря по тому, рассматривается ли $a(x)$ как базисный полином или как норма, --- распадается на два отдельных разложения: на представление главного идеала как наименьшего общего кратного наибольших первичных компонент по I, что здесь сводится к представлению как произведения степеней простых идеалов; или на разложение нормы как произведения степеней норм отдельных простых идеалов по IV. В построении поля нулей по пункту 2 на самом деле речь идет о поле нулей простого идеала, определенного $p(x)$, тогда как в пункте 3 на линейные множители разлагается норма. Определение кратности по пункту 4 подчиняется данному в IV, поскольку основной идеал нулевой ступени равен единичному идеалу, а основной идеал первой ступени при одной переменной уже равен самому идеалу.

\emph{Включение в обычную теорию элиминации} задается тем, что норма идеала имеет значение результанта; при таком понимании старая проблема теории элиминации становится частью теории идеалов полиномиальной области, а именно той частью, которая наряду с теоремами о разложении идеалов также рассматривает разложение их норм и связь обоих разложений.

\begin{center}
(Поступило 19. 10. 1923.)
\end{center}

\end{document}
